\documentclass[10pt,a4paper]{article}

\usepackage{fontspec}
\usepackage{polyglossia}
\setdefaultlanguage{english}
\IfFontExistsTF{STIXTwoText-Regular.otf}{
  \setmainfont{STIXTwoText}[
    Extension=.otf,
    UprightFont=*-Regular, ItalicFont=*-Italic,
    BoldFont=*-Bold, BoldItalicFont=*-BoldItalic
  ]
}{\setmainfont{PT Sans}}
\usepackage[a4paper,top=1.15cm,bottom=1.15cm,left=1.65cm,right=1.65cm]{geometry}
\usepackage{titlesec}
\usepackage{enumitem}
\usepackage[hidelinks]{hyperref}
\hypersetup{pdfauthor={Ilya Mikhailovskiy},pdftitle={Ilya Mikhailovskiy | Machine Learning Engineer | NLP and LLMs},pdflang=en-US}
\urlstyle{same}

\titleformat{\section}{\fontsize{11.7}{13.5}\selectfont}{}{0pt}{\MakeUppercase}[\vspace{0.6pt}{\titlerule[0.45pt]}]
\titlespacing*{\section}{0pt}{8pt}{4pt}
\setlist[itemize]{leftmargin=1.65em,labelsep=0.55em,itemsep=1.2pt,topsep=1pt,parsep=0pt,partopsep=0pt}
\renewcommand{\labelitemi}{\textbullet}
\newcommand{\role}[2]{\textbf{#1}\hfill{\small #2}\par}
\newcommand{\sep}{\enspace\textbar\enspace}
\setlength{\parindent}{0pt}
\setlength{\parskip}{0pt}
\pagestyle{empty}
\raggedright

\begin{document}
\fontsize{10.3}{12.2}\selectfont
{\centering
  {\fontsize{23}{26}\selectfont\bfseries Ilya Mikhailovskiy}\par
  \vspace{1.5pt}
  {\large Machine Learning Engineer | NLP and LLMs}\par
  \vspace{2pt}
  Moscow\sep On-site, hybrid or remote\sep English (C1)\par
  \href{mailto:ilyamihailovsy@gmail.com}{\underline{ilyamihailovsy@gmail.com}}\sep
  \href{https://t.me/Iluuuaaa}{\underline{Telegram @Iluuuaaa}}\sep
  \href{https://github.com/iluuua}{\underline{github.com/iluuua}}\par
}

\section{Experience}
\role{ReachFlow \normalfont\enspace\textit{Founder and Developer}\enspace\href{https://reachflow.tech}{\underline{reachflow.tech}}}{Jan 2026 to present}
Software for sales teams that work with customers on Telegram.
\begin{itemize}
  \item Built an LLM agent that plans replies and checks them against the conversation history. Cut classifier calls from five to two, reducing input text by 38\% in a representative test conversation.
  \item Built a prospecting tool that uses breadth-first search to find Telegram communities and an LLM to identify relevant messages. It adds qualified contacts and deals to the CRM without creating duplicates.
  \item Evaluated four LLMs on the same 35 real message drafts. Measured cost per call and checked whether each model preserved the approved wording when splitting a draft into shorter messages.
  \item Built a service that routes requests through OpenRouter by task, switches to a backup model on failure, validates responses with Pydantic and tracks the cost of each call.
\end{itemize}

\vspace{3pt}
\role{Delnovo \normalfont\enspace\textit{Chief Technology Officer}}{Jul 2023 to May 2026}
\begin{itemize}
  \item Built a support bot for a window distributor using retrieval-augmented generation and vector search over its PostgreSQL knowledge base. The client estimated that the bot handled about 80\% of routine inquiries.
  \item Connected a Telegram bot to the Bitrix24 REST API so customers could receive order status updates in chat.
  \item Built an OAuth-protected API for AI agents to search and read a client's catalog. Documented it with OpenAPI and automated publishing for 89 pages in Markdown for agents and HTML with JSON-LD for search engines.
\end{itemize}

\vspace{3pt}
\role{NEAR AI \normalfont\enspace\textit{Solutions Engineer}}{Jun 2023 to Jul 2024}
\begin{itemize}
  \item Solved 300 competitive programming problems in C++ for AI training, including 15 at International Olympiad in Informatics (IOI) level. Wrote editorials in LaTeX, analyzed algorithm complexity and reviewed solutions from teammates in other countries.
\end{itemize}

\section{Projects and competitions}
\role{Leaders of Digital Transformation Hackathon \normalfont\enspace\textit{Moscow Metro challenge}}{Oct 2026}
\begin{itemize}
  \item Led xixi-square, a two-person team that placed \href{https://i.moscow/upload/lending/lct2026/edebdf.pdf\#page=5}{\underline{second}}. We built a C++ obstacle detector for 3D lidar and integrated it with ROS~2. \href{https://github.com/iluuua/metro-lidar-obstacle-detection}{\underline{Code and demo}}
  \item Used the rails and tunnel cross section to estimate the train's path and check for obstacles up to 300~m ahead. The detector averaged 13.9~ms per frame on CPU and recorded zero false alarms over 15.4~km of clear track.
\end{itemize}

\vspace{3pt}
\role{Artificial Intelligence and Data Analysis Olympiad (AIDAO)}{Nov 2025}
\begin{itemize}
  \item Placed sixth out of 248 teams. Implemented \href{https://github.com/iluuua/aidao-2025-bev-occupancy}{\underline{Lift-Splat-Shoot in PyTorch}} from scratch to map static obstacles around a vehicle in bird's-eye view using images from four cameras.
  \item Recomputed the feature projection for each frame using camera calibration. Scored 0.5606 IoU on the private test set, compared with the winning team's 0.564.
\end{itemize}

\section{Skills}
\textbf{ML and NLP}\quad PyTorch, NumPy, pandas, ResNet, FPN, segmentation, RAG, embeddings, pgvector, JSON Schema, tool calling, LLM evaluation.\par
\textbf{Engineering}\quad Python, C++, FastAPI, asyncio, PostgreSQL, SQL, Redis, Docker, Linux, Git, CI/CD, pytest, Playwright.\par
\textbf{AI tools}\quad Claude, Codex, personal ChatGPT Pro subscription (\$200/month).

\section{Education and research}
\textbf{Moscow Polytechnic University}, Applied Mathematics and Computer Science, expected 2029. Auditing courses at HSE University's Faculty of Computer Science.\par
\textbf{Coursework (2025).} Deep Learning School at MIPT. January School in Mathematics and Programming at Sirius (92 academic hours). Algorithms, Track B, T-Bank's T-Generation program.\par
\textbf{Research.} Wrote a paper on stresses in aluminum around an Fe₄Al₁₃ inclusion. \href{https://github.com/iluuua/magnetostriction-plasticity-bounds}{\underline{Code and materials}}

\end{document}
